% Generated draft: requires mathematical and source review.
\documentclass[11pt]{article}
\usepackage[margin=25mm]{geometry}
\usepackage{fontspec}
\setmainfont{Latin Modern Roman}
\usepackage{amsmath,amssymb,mathtools,unicode-math}
\setmathfont{Latin Modern Math}
\usepackage{enumitem,needspace,xurl,hyperref}
\setlength{\parindent}{0pt}
\setlength{\parskip}{5pt}
\setlength{\emergencystretch}{3em}
\newcommand{\N}{\mathbb{N}}
\newcommand{\Z}{\mathbb{Z}}
\newcommand{\Q}{\mathbb{Q}}
\newcommand{\R}{\mathbb{R}}
\newcommand{\C}{\mathbb{C}}
\newcommand{\F}{\mathbb{F}}
\newcommand{\A}{\mathbb{A}}
\newcommand{\PP}{\mathbb{P}}
\newcommand{\E}{\mathbb{E}}
\newcommand{\eps}{\varepsilon}
\newcommand{\dd}{\,\mathrm{d}}
\begin{document}
\section*{UnsolvedMath ID: 30002878}
\textbf{Source problem code:} catalogue entry
\subsection{Definitions and mathematical statement}
Let \(X\) be a finite CW complex. A complex Azumaya algebra bundle of degree \(n\geq 1\) on \(X\) is a locally trivial bundle of unital complex algebras with fibre \(M_n(\C)\). Write \(\operatorname{Br}(X)\) for the topological Brauer group of Morita equivalence classes of such bundles, with group operation induced by tensor product. For \(\alpha\in\operatorname{Br}(X)\), define its topological period and index by
\[
\operatorname{per}_{\mathrm{topo}}(\alpha)=\operatorname{ord}_{\operatorname{Br}(X)}(\alpha),\qquad
\operatorname{ind}_{\mathrm{topo}}(\alpha)=\gcd\bigl\{n\geq 1: \text{an Azumaya algebra bundle of degree }n\text{ represents }\alpha\bigr\}.
\]
The period is finite, and the zero class has period and index \(1\). No connectedness assumption is supplied; in the degree condition above, the fibre is \(M_n(\C)\) throughout \(X\).

\textbf{Supplied conjectural formulation.} For every integer \(d\geq 1\), every finite CW complex \(X\) of dimension \(2d\), and every \(\alpha\in\operatorname{Br}(X)\), put \(m=\operatorname{per}_{\mathrm{topo}}(\alpha)\). The proposed assertion is
\[
\bigl(d\neq 3\ \text{or}\ m\not\equiv 2\pmod 4\bigr)
\quad\Longrightarrow\quad
\operatorname{ind}_{\mathrm{topo}}(\alpha)\mid m^{d-1}.
\]
The supplied statement makes no assertion for \(d=3\) and \(m\equiv 2\pmod 4\); its phrase ``except possibly'' does not assert failure for every class in that range or supply a replacement bound.

\textbf{Unresolved qualifications.} The source does not specify the range of \(d\); the display uses the positive-integer interpretation, and treatment of \(d=0\) requires clarification. The dimension condition is retained as \(\dim X=2d\), without replacing it by an upper bound. The source's \(\operatorname{per}(\alpha)\) in the exceptional condition is interpreted as the same topological period used elsewhere. The supplied dated curation calls this an unresolved all-dimensions ``straw-man'' formulation with partial results; this draft preserves that qualification and does not establish its present validity or status.
\subsection{Short English statement}
The supplied conjecture asks whether the topological index of every Brauer class on a finite CW complex of dimension \(2d\) divides the \((d-1)\)-st power of its period. It excludes from the assertion the possible exceptional case \(d=3\) with period congruent to \(2\) modulo \(4\). The supplied partially solved label and literature assessment require independent review.
\subsection{Sources}
\begin{itemize}
\item Supplied local UnsolvedMath registry record, permanent ID \(30002878\), entitled \emph{Topological Period--Index Conjecture}; exact source problem code \texttt{catalogue entry}, in the Oberwolfach Reports Open Problems collection. The record gives proposed year 2015 and names no individual proposer.
\item Supplied report citation: \emph{Cohomology of Finite Groups: Interactions and Applications} (2016), Oberwolfach Reports, DOI 10.4171/owr/2015/24, pp.\ 1317--1384. \url{https://doi.org/10.4171/owr/2015/24}. The background describes Oberwolfach Reports as documenting mathematical workshops and their open-problem discussions. The supplied 2016 citation year and 2015 proposed-year/DOI information are retained without resolving their relationship.
\item Antieau--Williams, \emph{The period-index problem for twisted topological K-theory}. Supplied literature reference: \url{https://doi.org/10.1112/jtopol/jtt042}.
\item Antieau--Williams, \emph{The topological period-index conjecture}. Supplied literature reference: \url{https://arxiv.org/abs/2003.10539}.
\item Literature assessment dated \texttt{2026-08-21}: the supplied curation labels the problem partially solved and states that the low-dimensional cases and the sharp period-two dimension-three exception are established, that general prime-dependent bounds are known, and that the exact all-dimensions straw-man formulation remains unresolved. This is supplied dated curation derived from the OWR record, not an independently verified literature assessment; its time-sensitive status and cited claims require independent verification.
\item The registry labels the statement \texttt{corrected\_verified}, explaining that mathematical content was preserved while notation, typography, or omitted definitions were repaired for publication. This is a supplied editorial label, not mathematical verification performed for this draft.
\end{itemize}
\end{document}
